% ============================================================
%  MIPSolvers 使用教程 — Beamer 幻灯片
%  编译: xelatex -interaction=nonstopmode slides.tex  (执行两次)
% ============================================================
\documentclass[10pt,aspectratio=169]{beamer}

%% ── 中文支持 ──────────────────────────────────────────────────────────────────
\usepackage{ctex}

%% ── 主题 ──────────────────────────────────────────────────────────────────────
\usetheme{Madrid}
\usecolortheme{seahorse}
\setbeamertemplate{navigation symbols}{}
\setbeamertemplate{footline}[frame number]

%% ── 数学 ──────────────────────────────────────────────────────────────────────
\usepackage{amsmath,amssymb,bm}
\newcommand{\vx}{\bm{x}}
\newcommand{\vc}{\bm{c}}
\newcommand{\vb}{\bm{b}}
\newcommand{\mA}{\bm{A}}
\newcommand{\mQ}{\bm{Q}}
\newcommand{\R}{\mathbb{R}}
\newcommand{\Z}{\mathbb{Z}}

%% ── 表格 ──────────────────────────────────────────────────────────────────────
\usepackage{booktabs,multirow,array}

%% ── 代码高亮 ──────────────────────────────────────────────────────────────────
\usepackage{listings,xcolor}

\definecolor{codebg}  {RGB}{248,248,248}
\definecolor{codecmmt}{RGB}{ 60,128, 49}
\definecolor{codekey} {RGB}{ 22, 22,212}
\definecolor{codestr} {RGB}{163,  21, 21}
\definecolor{codegray}{RGB}{130,130,130}

\lstdefinestyle{py}{
  backgroundcolor = \color{codebg},
  basicstyle      = \ttfamily\scriptsize,
  keywordstyle    = \color{codekey}\bfseries,
  commentstyle    = \color{codecmmt}\itshape,
  stringstyle     = \color{codestr},
  numberstyle     = \tiny\color{codegray},
  language        = Python,
  numbers         = left,
  numbersep       = 5pt,
  stepnumber      = 1,
  showstringspaces= false,
  breaklines      = true,
  frame           = single,
  rulecolor       = \color{gray!40},
  tabsize         = 4,
  keepspaces      = true,
  columns         = fullflexible,
}
\lstset{style=py}

%% ── 颜色方框 ─────────────────────────────────────────────────────────────────
\usepackage{tcolorbox}
\tcbuselibrary{skins,breakable}
\tcbset{
  mybox/.style={
    colback=blue!5,colframe=blue!60!black,
    fonttitle=\bfseries,boxrule=0.5pt,
    left=2pt,right=2pt,top=2pt,bottom=2pt
  }
}

%% ── TikZ 流程图 ──────────────────────────────────────────────────────────────
\usepackage{tikz}
\usetikzlibrary{arrows.meta,shapes.geometric,positioning,fit}

%% ── 标题信息 ─────────────────────────────────────────────────────────────────
\title{MIPSolvers 使用教程}
\subtitle{优化求解器：理论、接口与电力系统应用}
\author{内部技术教程}
\institute{MIPSolvers v0.1}
\date{\today}

% ============================================================
\begin{document}
% ============================================================

%% ── 封面 ──────────────────────────────────────────────────────────────────────
\begin{frame}
  \titlepage
  \vfill
  \begin{center}
    \small\texttt{Python~3.8 + pybind11 + HiGHS/SCIP}
  \end{center}
\end{frame}

%% ── 目录 ──────────────────────────────────────────────────────────────────────
\begin{frame}{目录}
  \tableofcontents
\end{frame}

% ============================================================
\section{优化求解器概论}
% ============================================================

\begin{frame}{什么是数学优化？}
  \begin{columns}[T]
    \column{0.55\textwidth}
    \textbf{数学优化（Mathematical Programming）}
    \vspace{4pt}

    在满足约束条件的前提下，寻找使目标函数最小（或最大）的决策变量取值：
    \[
      \min_{\vx\in\mathcal{X}} \; f(\vx)
      \quad\text{s.t.}\quad
      g_i(\vx) \le 0,\;
      h_j(\vx) = 0
    \]

    \textbf{三要素：}
    \begin{itemize}
      \item \textbf{决策变量} $\vx$：可调节量
      \item \textbf{目标函数} $f(\vx)$：优化指标
      \item \textbf{约束条件} $g,\,h$：物理/逻辑限制
    \end{itemize}

    \column{0.40\textwidth}
    \begin{tcolorbox}[mybox,title=常见问题类型]
      \footnotesize
      \begin{tabular}{@{}ll@{}}
        \textbf{LP}   & 线性规划 \\[2pt]
        \textbf{QP}   & 二次规划 \\[2pt]
        \textbf{NLP}  & 非线性规划 \\[2pt]
        \textbf{MILP} & 混合整数线性规划 \\[2pt]
        \textbf{MINLP}& 混合整数非线性 \\[2pt]
        \textbf{LE}   & 线性方程组 \\[2pt]
        \textbf{NLE}  & 非线性方程组 \\
      \end{tabular}
    \end{tcolorbox}
  \end{columns}
\end{frame}

\begin{frame}{求解器的作用与分类}
  \begin{columns}[T]
    \column{0.5\textwidth}
    \textbf{求解器}封装了数值优化算法，提供统一接口：
    \begin{center}
      \begin{tikzpicture}[>=Stealth,node distance=0.9cm,font=\small]
        \node[draw,rounded corners,fill=blue!15,text width=2.8cm,align=center]
              (model) {数学模型\\(变量, 目标, 约束)};
        \node[draw,rounded corners,fill=orange!20,text width=2.8cm,align=center,
              below=of model]
              (solver){优化求解器\\(算法实现)};
        \node[draw,rounded corners,fill=green!20,text width=2.8cm,align=center,
              below=of solver]
              (sol)  {最优解 / 状态};
        \draw[->] (model)  -- (solver);
        \draw[->] (solver) -- (sol);
      \end{tikzpicture}
    \end{center}

    \column{0.45\textwidth}
    \footnotesize
    \begin{tabular}{@{}lll@{}}
      \toprule
      \textbf{求解器} & \textbf{类型} & \textbf{许可} \\
      \midrule
      Gurobi      & LP/MILP  & 商业 \\
      CPLEX       & LP/MILP  & 商业 \\
      HiGHS       & LP/MILP  & 开源 \\
      SCIP        & MILP     & 开源 \\
      Ipopt       & NLP      & 开源 \\
      GLPK        & LP/MILP  & 开源 \\
      \midrule
      \multicolumn{3}{l}{\textit{MIPSolvers 统一以上接口}} \\
      \bottomrule
    \end{tabular}
    \vspace{6pt}

    \textbf{核心算法：}
    \begin{itemize}
      \item 单纯形法 / 对偶单纯形法 (LP)
      \item 内点法 (LP, NLP)
      \item 分支定界法 (MILP)
      \item 切平面法、列生成法
    \end{itemize}
  \end{columns}
\end{frame}

\begin{frame}{主要应用场景}
  \begin{columns}[T]
    \column{0.48\textwidth}
    \textbf{通用领域：}
    \begin{itemize}
      \item 供应链与物流优化（运输规划、选址）
      \item 金融投资组合优化
      \item 生产排程与资源分配
      \item 机器学习（正则化、SVM）
      \item 通信网络路由
    \end{itemize}

    \vspace{8pt}
    \textbf{电力系统专项：}
    \begin{itemize}
      \item \textbf{机组组合（UC/SCUC）}：确定每台机组的开/停计划
      \item \textbf{最优潮流（OPF）}：最小化发电成本，满足电网约束
      \item \textbf{电力市场出清}：计算边际电价（LMP）
      \item \textbf{输电规划}：确定新建线路方案
      \item \textbf{储能调度}：管理电池充放电计划
    \end{itemize}

    \column{0.47\textwidth}
    \begin{tcolorbox}[mybox,title=电力系统调度层级]
      \footnotesize
      \begin{tabular}{@{}ll@{}}
        \textbf{年度/季度} & 检修计划规划 \\[2pt]
        \textbf{周}        & 水电协调优化 \\[2pt]
        \textbf{日前}      & SCUC（机组组合）\\[2pt]
        \textbf{小时前}    & SCED（经济调度）\\[2pt]
        \textbf{实时}      & AGC自动发电控制 \\[2pt]
        \textbf{毫秒级}    & 保护与稳定控制 \\
      \end{tabular}
    \end{tcolorbox}

    \vspace{6pt}
    \textbf{MIPSolvers 重点}：日前 SCUC 及 SCED 的快速求解与建模
  \end{columns}
\end{frame}

\begin{frame}{电力系统优化的理论难题（一）}
  \begin{columns}[T]
    \column{0.52\textwidth}
    \textbf{1. 计算复杂性（NP难）}

    SCUC 中有 $G$ 台机组、$T$ 个时段：
    \[
      \text{组合数} = 2^{G \times T}
    \]
    \begin{itemize}
      \item $G=50,\; T=24$：超过 $10^{360}$ 种方案！
      \item 启停约束（最小开/停机时间）引入时间耦合
      \item 分段报价将连续变量线性化，增加辅助二值变量
    \end{itemize}

    \vspace{6pt}
    \textbf{2. 时间耦合约束}
    \begin{itemize}
      \item 爬坡约束：$|p_t - p_{t-1}| \le R_g$
      \item 最小开机时间：$\sum_{\tau=t}^{t+T^{up}-1} u_\tau \ge T^{up} \cdot (u_t - u_{t-1})$
      \item 使 LP 松弛的质量决定 B\&B 效率
    \end{itemize}

    \column{0.43\textwidth}
    \begin{tcolorbox}[mybox,title=实际规模（华东电网示例）]
      \footnotesize
      \begin{tabular}{@{}lr@{}}
        机组数      & $\sim$1000 \\
        节点数      & $\sim$2000 \\
        线路数      & $\sim$3000 \\
        时段数      & 24--96 \\
        二值变量    & $> 24000$ \\
        求解时限    & $\le 30$~min \\
      \end{tabular}
    \end{tcolorbox}

    \vspace{8pt}
    \textbf{常用应对策略：}
    \begin{itemize}
      \item Lagrangian 松弛
      \item Benders 分解
      \item 强化形式（Tight formulations）
      \item 列生成 / 分支定价
      \item 热启动（Warm start）
    \end{itemize}
  \end{columns}
\end{frame}

\begin{frame}{电力系统优化的理论难题（二）}
  \begin{columns}[T]
    \column{0.50\textwidth}
    \textbf{3. 非凸性（AC-OPF）}

    交流最优潮流引入非线性等式约束：
    \[
      P_i = \sum_j |V_i||V_j|\left(G_{ij}\cos\theta_{ij} + B_{ij}\sin\theta_{ij}\right)
    \]
    \begin{itemize}
      \item 电压模/相角乘积为非凸
      \item 局部最优解众多，全局优化困难
      \item SDP / 多项式优化松弛方法
    \end{itemize}

    \vspace{6pt}
    \textbf{4. 不确定性（随机优化）}
    \begin{itemize}
      \item 风电/光伏出力不确定
      \item 负荷预测误差
      \item 随机 UC、鲁棒 UC、机会约束
    \end{itemize}

    \column{0.45\textwidth}
    \textbf{5. 数值挑战}
    \begin{itemize}
      \item Big-M 取值不当→数值病态
      \item 线路潮流约束引入大系数差异
      \item 稀疏 vs. 密集矩阵混合
    \end{itemize}

    \vspace{10pt}
    \begin{tcolorbox}[mybox,title=MIPSolvers 目标]
      \footnotesize
      提供\textbf{统一的求解器接口}，使研究人员专注于：
      \begin{itemize}
        \item 建模与算法设计
        \item 快速原型验证
        \item 多求解器对比
      \end{itemize}
      而不是底层求解器调用的细节
    \end{tcolorbox}
  \end{columns}
\end{frame}

% ============================================================
\section{Engine 数学格式}
% ============================================================

\begin{frame}{支持的标准数学格式总览}
  \begin{center}
  \footnotesize
  \begin{tabular}{@{}lllp{5.2cm}@{}}
    \toprule
    \textbf{类型} & \textbf{目标} & \textbf{变量域} & \textbf{可用求解器} \\
    \midrule
    LP    & 线性    & 连续   & HiGHS, Gurobi, NativeLCQP, NativePDLP, NativeIPMLP \\[2pt]
    QP    & 二次    & 连续   & HiGHS, Gurobi, NativeLCQP \\[2pt]
    NLP   & 非线性  & 连续   & NativeIPMLP（内点法） \\[2pt]
    MILP  & 线性    & 混合整数 & HiGHS, Gurobi, SCIP, NativeBranchAndCut \\[2pt]
    MINLP & 非线性  & 混合整数 & NativeBranchAndCut \\[2pt]
    LE    & ---     & 连续   & NativeDirect（稀疏 LU 分解）\\[2pt]
    NLE   & ---     & 连续   & NativeNewton（Newton-Raphson）\\
    \bottomrule
  \end{tabular}
  \end{center}

  \vspace{8pt}
  \textbf{统一入口}：\texttt{mipsolvers.engine.solve\_lp / solve\_milp / solve\_nlp / ...}

  \vspace{4pt}
  \textbf{自动选择}：\texttt{solver="Auto"} 按优先级选取已安装且最适合的求解器
\end{frame}

\begin{frame}{LP — 线性规划}
  \begin{columns}[T]
    \column{0.50\textwidth}
    \textbf{标准形式：}
    \[
      \min_{\vx}\; \vc^{\top}\vx
    \]
    \[
      \text{s.t.}\quad
      \begin{cases}
        \mA\,\vx \le \vb \\
        A_{eq}\vx = b_{eq} \\
        l_b \le \vx \le u_b
      \end{cases}
    \]

    \textbf{Python 接口参数：}
    \begin{itemize}
      \item \texttt{c}：目标系数向量 $\R^n$
      \item \texttt{A, b}：不等式约束
      \item \texttt{Aeq, beq}：等式约束（可省略）
      \item \texttt{lb, ub}：变量界（可省略）
      \item \texttt{solver}：求解器名称
      \item \texttt{maximize}：默认\,\texttt{False}
    \end{itemize}

    \column{0.46\textwidth}
    \textbf{返回字典 \texttt{res}：}
    \begin{itemize}
      \item \texttt{res['success']}：是否求解成功
      \item \texttt{res['objective']}：最优目标值
      \item \texttt{res['x']}：最优解向量
      \item \texttt{res['status']}：状态描述
      \item \texttt{res['duals']}：对偶变量（影子价格）
      \item \texttt{res['solver']}：实际使用求解器
      \item \texttt{res['runtime\_sec']}：耗时
    \end{itemize}

    \vspace{6pt}
    \begin{tcolorbox}[mybox,title=应用示例]
      \footnotesize
      生产计划、运输、网络流、电力经济调度
    \end{tcolorbox}
  \end{columns}
\end{frame}

\begin{frame}{MILP — 混合整数线性规划}
  \begin{columns}[T]
    \column{0.50\textwidth}
    \textbf{标准形式：}
    \[
      \min_{\vx}\; \vc^{\top}\vx
      \quad\text{s.t.}\quad
      \mA\vx\le\vb,\; l_b\le\vx\le u_b
    \]
    \[
      x_i \in \mathbb{Z} \;(\text{或}\; \{0,1\}) \quad \forall i \in \mathcal{I}
    \]

    \textbf{额外参数：}
    \begin{itemize}
      \item \texttt{vartypes}：列表，每元素为\\
            \texttt{"C"}（连续）、\texttt{"B"}（0/1）、\texttt{"I"}（整数）
      \item \texttt{mip\_gap}：MIP容忍间隙（默认 $10^{-4}$）
      \item \texttt{time\_limit\_sec}：时间上限
    \end{itemize}

    \column{0.45\textwidth}
    \textbf{额外返回字段：}
    \begin{itemize}
      \item \texttt{res['mip\_gap']}：实际间隙
    \end{itemize}

    \vspace{8pt}
    \textbf{算法：Branch \& Bound}
    \begin{enumerate}
      \item 求解 LP 松弛，得到下界
      \item 选择分数变量进行分支
      \item 剪枝：界+LP松弛判断
      \item 更新上界（可行整数解）
    \end{enumerate}

    \vspace{6pt}
    \begin{tcolorbox}[mybox,title=应用示例]
      \footnotesize
      背包/排程、设施选址、\\
      机组组合、供应链规划
    \end{tcolorbox}
  \end{columns}
\end{frame}

\begin{frame}{QP、NLP 与方程组求解}
  \begin{columns}[T]
    \column{0.48\textwidth}
    \textbf{QP — 二次规划：}
    \[
      \min\; \tfrac{1}{2}\vx^{\top}\mQ\vx + \vc^{\top}\vx
      \quad\text{s.t.}\quad \mA\vx\le\vb
    \]
    \begin{itemize}
      \item $\mQ \succeq 0$（正半定）→ 凸问题
      \item 参数：\texttt{Q}（稀疏/密集），其余同 LP
      \item 应用：投资组合、MPC 轨迹优化
    \end{itemize}

    \vspace{6pt}
    \textbf{NLP — 非线性规划：}
    \begin{itemize}
      \item 用户提供 $f$, $\nabla f$, $H_f$, $g$, $\nabla g$ 回调函数
      \item 内点法 (NativeIPMLP)
      \item 应用：AC-OPF、非线性工程优化
    \end{itemize}

    \column{0.48\textwidth}
    \textbf{LE — 线性方程组：}
    \[
      \mA\vx = \vb,\quad \mA\in\R^{m\times n}
    \]
    \begin{itemize}
      \item 稀疏 LU 分解（NativeDirect）
      \item 适用：潮流方程线性化、预条件
    \end{itemize}

    \vspace{6pt}
    \textbf{NLE — 非线性方程组：}
    \[
      F(\vx) = \bm{0},\quad F:\R^n\to\R^n
    \]
    \begin{itemize}
      \item Newton-Raphson 方法
      \item 需提供 $F$ 和 $J_F$
      \item 适用：潮流计算、稳态方程
    \end{itemize}
  \end{columns}
\end{frame}

% ============================================================
\section{Python 接口使用}
% ============================================================

\begin{frame}[fragile]{两种 Python 接口}
  \begin{columns}[T]
    \column{0.48\textwidth}
    \begin{tcolorbox}[colback=orange!8,colframe=orange!70,title=底层 Engine 接口]
      \footnotesize
      \begin{itemize}
        \item 直接传 NumPy 矩阵/向量
        \item 适合矩阵结构已知的小规模问题
        \item 最低开销，最直接控制
        \item 返回字典（\texttt{res['x']}, \texttt{res['objective']}）
      \end{itemize}
      \vspace{4pt}
      \texttt{mipsolvers.engine.solve\_lp(...)}\\
      \texttt{mipsolvers.engine.solve\_milp(...)}
    \end{tcolorbox}

    \column{0.48\textwidth}
    \begin{tcolorbox}[colback=green!8,colframe=green!60!black,title=AML 代数建模接口]
      \footnotesize
      \begin{itemize}
        \item 声明式建模：集合、参数、变量、约束
        \item 支持索引变量（多维下标）
        \item 可读性强，接近数学符号
        \item 自动管理矩阵构造
        \item 支持 LP / MILP / QP / NLP
      \end{itemize}
      \vspace{4pt}
      \texttt{mipsolvers.aml.Model(...)}
    \end{tcolorbox}
  \end{columns}

  \vspace{10pt}
  \begin{center}
  \textbf{导入模块：}
  \end{center}
  \vspace{-4pt}
  \begin{columns}
    \column{0.5\textwidth}
    \begin{lstlisting}
import sys
sys.path.insert(0,
    '/path/to/MIPSolvers/build_py')
import mipsolvers
aml = mipsolvers.aml
    \end{lstlisting}
  \end{columns}
\end{frame}

\begin{frame}[fragile]{底层接口：求解 LP（生产计划）}
  \begin{columns}[T]
    \column{0.52\textwidth}
    \textbf{问题：} $\max\;5x_1+4x_2$
    \[
      \text{s.t.}\;\;
      \begin{pmatrix}6&4\\1&2\end{pmatrix}
      \begin{pmatrix}x_1\\x_2\end{pmatrix}
      \le\begin{pmatrix}24\\6\end{pmatrix},\;
      x_1,x_2\ge0
    \]

    \begin{lstlisting}
import numpy as np
import mipsolvers

c  = np.array([-5.0, -4.0])  # min
A  = np.array([[6.0,4.0],[1.0,2.0]])
b  = np.array([24.0, 6.0])
lb = np.zeros(2)

res = mipsolvers.engine.solve_lp(
    c=c, A=A, b=b, lb=lb,
    solver="Auto")
    \end{lstlisting}

    \column{0.44\textwidth}
    \begin{lstlisting}
print(res['status'])
# Optimal

print(-res['objective'])
# 21.0

print(res['x'])
# [3.0  1.5]

print(res['duals'])
# [0.75  0.5]  (shadow prices)
    \end{lstlisting}

    \vspace{8pt}
    \textbf{完整示例：} \texttt{demo\_01\_lp\_direct.py}

    \vspace{4pt}
    \begin{tcolorbox}[mybox,title=查询可用求解器]
      \footnotesize
      \texttt{mipsolvers.engine.list\_solvers("LP")}
    \end{tcolorbox}
  \end{columns}
\end{frame}

\begin{frame}[fragile]{底层接口：求解 MILP（0/1 背包）}
  \begin{columns}[T]
    \column{0.52\textwidth}
    \textbf{问题：} 选物品使总价值最大，总重量 $\le W$
    \[
      \max\;\bm{v}^{\top}\vx,\quad
      \bm{w}^{\top}\vx\le W,\quad
      x_i\in\{0,1\}
    \]

    \begin{lstlisting}
weights = np.array([3,4,5,2,6,1],
                   dtype=float)
values  = np.array([9,8,10,5,12,3],
                   dtype=float)
W, n = 10.0, 6

res = mipsolvers.engine.solve_milp(
    c   = -values,
    A   = weights.reshape(1, n),
    b   = np.array([W]),
    lb  = np.zeros(n),
    ub  = np.ones(n),
    vartypes = ["B"] * n,
    solver   = "Auto")
    \end{lstlisting}

    \column{0.44\textwidth}
    \begin{lstlisting}
print(-res['objective'])
# 22.0

sel = [i for i in range(n)
       if res['x'][i] > 0.5]
# 选了哪些物品

print(res['mip_gap'])
# 0.0  (精确最优)

print(res['solver'])
# HiGHS
    \end{lstlisting}

    \vspace{6pt}
    \textbf{完整示例：} \texttt{demo\_02\_milp\_knapsack.py}

    \vspace{4pt}
    \begin{tcolorbox}[mybox,title=变量类型]
      \footnotesize
      \texttt{"B"} = 0/1 整型\\
      \texttt{"I"} = 普通整型\\
      \texttt{"C"} = 连续型
    \end{tcolorbox}
  \end{columns}
\end{frame}

\begin{frame}{AML 代数建模工作流程}
  \begin{center}
  \begin{tikzpicture}[>=Stealth,node distance=0.45cm,font=\small]
    \tikzstyle{box} = [draw,rounded corners,fill=blue!12,
                       text width=2.2cm,align=center,
                       minimum height=0.7cm]
    \tikzstyle{arr} = [->,thick,blue!60]

    \node[box] (m)  {创建模型\\{\tiny\texttt{aml.Model("name")}}};
    \node[box, right=0.7cm of m] (s)  {定义集合\\{\tiny\texttt{add\_set / add\_ordered\_set}}};
    \node[box, right=0.7cm of s] (p)  {载入参数\\{\tiny\texttt{add\_param / load()}}};
    \node[box, right=0.7cm of p] (v)  {声明变量\\{\tiny\texttt{add\_var2}}};
    \node[box, right=0.7cm of v] (oc) {目标+约束\\{\tiny\texttt{minimize / add\_constraints}}};
    \node[box, right=0.7cm of oc](sl) {求解\\{\tiny\texttt{model.solve(opts)}}};

    \draw[arr] (m) -- (s);
    \draw[arr] (s) -- (p);
    \draw[arr] (p) -- (v);
    \draw[arr] (v) -- (oc);
    \draw[arr] (oc)-- (sl);

    \node[below=0.5cm of sl, align=center, font=\footnotesize]
      {$\downarrow$\\
       \texttt{SolveResult}\\
       {\tiny \texttt{termination\_status, objective\_value}}\\
       {\tiny \texttt{var\_value(), array\_values(), dual\_array()}}};
  \end{tikzpicture}
  \end{center}
\end{frame}

\begin{frame}[fragile]{AML：集合、参数与变量}
  \begin{columns}[T]
    \column{0.50\textwidth}
    \textbf{集合（索引域）：}
    \begin{lstlisting}
m = aml.Model("transport")

# 普通集合
I = m.add_set("I", ["A","B","C"])

# 有序集合（支持 prev/next）
T = m.add_ordered_set("T",
        ["t1","t2","t3","t4"])
    \end{lstlisting}

    \vspace{4pt}
    \textbf{参数：}
    \begin{lstlisting}
c = m.add_param("cost", dim=2)
c.load({
    ("A","t1"): 2.5,
    ("A","t2"): 3.0,
    ("B","t1"): 4.0,
    ...
})
val = c[("A","t1")]  # => 2.5
    \end{lstlisting}

    \column{0.46\textwidth}
    \textbf{决策变量：}
    \begin{lstlisting}
# 1维变量 (indexed by I)
x1 = m.add_var("x", I,
      aml.VarType.Continuous,
      lb=0.0)

# 2维变量 (indexed by I x T)
x2 = m.add_var2("y", I, T,
      aml.VarType.Binary)

# 访问方式
xv = x2[("A", "t1")]  # VarRef
    \end{lstlisting}

    \vspace{6pt}
    \textbf{变量类型：}
    \begin{itemize}
      \footnotesize
      \item \texttt{VarType.Continuous}
      \item \texttt{VarType.Binary}
      \item \texttt{VarType.Integer}
    \end{itemize}
  \end{columns}
\end{frame}

\begin{frame}[fragile]{AML：目标函数与约束}
  \begin{columns}[T]
    \column{0.52\textwidth}
    \textbf{目标函数（使用 sum\_over）：}
    \begin{lstlisting}
# min sum_{i,t} cost[i,t] * x[i,t]
m.minimize(
  aml.sum_over(I, lambda i:
    aml.sum_over(T, lambda t:
      cost[(i,t)] * x2[(i,t)])))
    \end{lstlisting}

    \vspace{4pt}
    \textbf{批量约束 add\_constraints：}
    \begin{lstlisting}
# 对集合I中每个i添加约束
m.add_constraints("cap", I,
  lambda i:
    aml.sum_over(T, lambda t:
      x2[(i,t)] * 1.0) <= cap[i])
    \end{lstlisting}

    \column{0.44\textwidth}
    \textbf{单个约束 add\_constraint：}
    \begin{lstlisting}
# 线性表达式比较
lhs = 1.0*x1[i] - 50.0*x2[i]
m.add_constraint(
    lhs >= 0.0, "pmin_i")

# 表达式 == 常数
m.add_constraint(
    1.0*a + 1.0*b == demand,
    "balance")
    \end{lstlisting}

    \vspace{6pt}
    \begin{tcolorbox}[mybox,title=表达式构造规则]
      \footnotesize
      \texttt{float * VarRef} $\to$ \texttt{LinearExpr}\\
      \texttt{LinearExpr $\pm$ LinearExpr} $\to$ \texttt{LinearExpr}\\
      \texttt{LinearExpr $\ge$ float} $\to$ \texttt{TempConstr}
    \end{tcolorbox}
  \end{columns}
\end{frame}

\begin{frame}[fragile]{AML：求解配置与结果提取}
  \begin{columns}[T]
    \column{0.50\textwidth}
    \textbf{求解选项：}
    \begin{lstlisting}
opts = aml.SolveOptions()
opts.solver_name  = ""    # Auto
opts.verbosity    = 0     # 0=silent
opts.mip_gap_tol  = 1e-3  # MIP gap
opts.time_limit   = 300.0 # seconds
opts.max_iter     = 10000

result = m.solve(opts)
    \end{lstlisting}

    \vspace{4pt}
    \textbf{基本结果：}
    \begin{lstlisting}
print(result.termination_status)
# TerminationStatus.Optimal

print(result.objective_value)
# 320.5

print(result.solver_used)
# HiGHS
    \end{lstlisting}

    \column{0.46\textwidth}
    \textbf{提取变量值：}
    \begin{lstlisting}
# 单个变量
v = result.var_value(x2[("A","t1")])

# 整个变量数组 -> dict
vals = result.array_values(x2)
# {("A","t1"): 2.5,
#  ("A","t2"): 0.0, ...}
for key, val in vals.items():
    if val > 0.01:
        print(key, val)
    \end{lstlisting}

    \vspace{4pt}
    \textbf{提取对偶变量：}
    \begin{lstlisting}
# 批量约束的对偶值
ca = m.add_constraints(
        "c", I, lambda i: ...)
duals = result.dual_array(ca)
# {"A": -0.75, "B": -0.50}
    \end{lstlisting}
  \end{columns}
\end{frame}

\begin{frame}[fragile]{AML 示例：运输问题（完整）}
  \begin{columns}[T]
    \column{0.52\textwidth}
    \begin{lstlisting}
m = aml.Model("transport")
F = m.add_set("F", ["FA","FB"])
C = m.add_set("C", ["C1","C2","C3"])
supply = m.add_param("s", dim=1)
demand = m.add_param("d", dim=1)
cost   = m.add_param("c", dim=2)
supply.load({"FA":120.0,"FB":80.0})
demand.load({"C1":70.0,"C2":90.0,
             "C3":40.0})
cost.load({("FA","C1"):4.0,
           ("FA","C2"):3.0,...})

x = m.add_var2("x", F, C,
    aml.VarType.Continuous, lb=0.0)

m.minimize(aml.sum_over(F, lambda f:
  aml.sum_over(C, lambda c:
    cost[(f,c)] * x[(f,c)])))
    \end{lstlisting}

    \column{0.44\textwidth}
    \begin{lstlisting}
m.add_constraints("sup", F,
  lambda f:
    aml.sum_over(C, lambda c:
      x[(f,c)]*1.0) <= supply[f])

m.add_constraints("dem", C,
  lambda c:
    aml.sum_over(F, lambda f:
      x[(f,c)]*1.0) >= demand[c])

opts = aml.SolveOptions()
r = m.solve(opts)
print(r.objective_value) # 570.0
    \end{lstlisting}

    \vspace{6pt}
    \textbf{完整示例：} \texttt{demo\_03\_aml\_transport.py}
  \end{columns}
\end{frame}

% ============================================================
\section{SCUC 问题建模与求解}
% ============================================================

\begin{frame}{SCUC 问题定义}
  \begin{columns}[T]
    \column{0.55\textwidth}
    \textbf{安全约束机组组合（SCUC）}

    \vspace{4pt}
    \textit{Security-Constrained Unit Commitment}

    \vspace{8pt}
    \textbf{本质}：日前调度问题，决定未来 $T=24$（或 96）个时段内：
    \begin{enumerate}
      \item 每台机组\textbf{何时开机/停机}（整数决策）
      \item 每台机组\textbf{出力多少 MW}（连续决策）
    \end{enumerate}

    \vspace{8pt}
    \textbf{优化目标}：最小化总运行成本（含启动费、空载费、发电费）

    \vspace{8pt}
    \textbf{约束体系}：
    \begin{itemize}
      \item 功率平衡（需量=供给）
      \item 出力上下限（取决于开停状态）
      \item 爬坡速率（相邻时段出力变化）
      \item 最小开/停机时间
      \item 网络安全约束（N-1 准则）
    \end{itemize}

    \column{0.40\textwidth}
    \begin{tcolorbox}[mybox,title=决策变量集合]
      \footnotesize
      \begin{tabular}{@{}ll@{}}
        $u_{g,t}\in\{0,1\}$   & 开停状态 \\[3pt]
        $p_{g,t}\ge 0$         & 出力 (MW) \\[3pt]
        $su_{g,t}\in\{0,1\}$  & 启动标志 \\[3pt]
        $sd_{g,t}\in\{0,1\}$  & 停机标志 \\[3pt]
        $r_{g,t}\ge 0$         & 备用量 (MW) \\
      \end{tabular}
    \end{tcolorbox}

    \vspace{8pt}
    \begin{tcolorbox}[colback=red!5,colframe=red!60,title=计算复杂度]
      \footnotesize
      SCUC 是 \textbf{NP难}问题\\
      $G$ 台机组，$T$ 时段：\\
      二值变量数 $= 3 \times G \times T$\\
      连续变量数 $= G \times T$
    \end{tcolorbox}
  \end{columns}
\end{frame}

\begin{frame}{SCUC 数学模型}
  \begin{columns}[T]
    \column{0.52\textwidth}
    \textbf{目标函数（最小化总成本）：}
    \[
      \min\sum_{g}\sum_{t}
      \Bigl[
        c^{nl}_g u_{g,t}
        + c^e_g p_{g,t}
        + c^{su}_g su_{g,t}
      \Bigr]
    \]

    其中 $c^{nl}$：空载成本，$c^e$：发电边际成本，$c^{su}$：启动成本

    \vspace{8pt}
    \textbf{功率平衡约束：}
    \[
      \sum_g p_{g,t} = D_t \quad \forall t
    \]
    （$D_t$：$t$时段系统负荷）

    \column{0.44\textwidth}
    \textbf{出力界约束（依赖开停状态）：}
    \[
      \underline{P}_g u_{g,t}
      \le p_{g,t}
      \le \overline{P}_g u_{g,t}
      \quad\forall g,t
    \]

    \textbf{爬坡约束：}
    \[
      p_{g,t} - p_{g,t-1} \le R_g \quad\forall g,t>1
    \]
    \[
      p_{g,t-1} - p_{g,t} \le R_g \quad\forall g,t>1
    \]

    \textbf{启停逻辑：}
    \[
      su_{g,t} - sd_{g,t}
      = u_{g,t} - u_{g,t-1}
      \quad\forall g,t
    \]
    \[
      su_{g,t} + sd_{g,t} \le 1
    \]
  \end{columns}
\end{frame}

\begin{frame}{SCUC 约束：最小开/停机时间}
  \begin{columns}[T]
    \column{0.55\textwidth}
    \textbf{最小开机时间（$T_g^{up}$）}：

    机组$g$在$t$期启动后，必须连续开机至少 $T_g^{up}$ 个时段：
    \[
      \sum_{\tau=t}^{\min(t+T^{up}_g-1,\,T)}
      u_{g,\tau}
      \ge T^{up}_g \cdot su_{g,t}
      \quad \forall t
    \]

    \vspace{8pt}
    \textbf{最小停机时间（$T_g^{dn}$）}：

    机组$g$在$t$期停机后，必须连续停机至少 $T_g^{dn}$ 个时段：
    \[
      \sum_{\tau=t}^{\min(t+T^{dn}_g-1,\,T)}
      (1-u_{g,\tau})
      \ge T^{dn}_g \cdot sd_{g,t}
      \quad \forall t
    \]

    \column{0.40\textwidth}
    \textbf{分段报价线性化：}

    机组出力成本用分段线性函数近似：
    \[
      c_g(p) = \sum_{k=1}^{K} c_{g,k}\,\delta_{g,k}
    \]
    其中 $\delta_{g,k}$ 为各段出力量，满足：
    \[
      \sum_k \delta_{g,k} = p_g,\quad
      0\le\delta_{g,k}\le\Delta\overline{P}_{g,k}
    \]

    \vspace{6pt}
    \begin{tcolorbox}[mybox,title=整体规模]
      \footnotesize
      典型日前调度（$G=100,T=24$）：\\
      $\sim 7200$ 二值变量\\
      $\sim 2400$ 连续变量\\
      $\sim 20000$ 约束\\
      求解时限：$\le 10$ 分钟
    \end{tcolorbox}
  \end{columns}
\end{frame}

\begin{frame}[fragile]{AML 建模 SCUC：数据与变量定义}
  \begin{columns}[T]
    \column{0.52\textwidth}
    \begin{lstlisting}
gen_data = {
  "G1": dict(pmin=50, pmax=200,
             ramp=120, no_load=500,
             bid=30, startup=1000),
  "G2": dict(pmin=30, pmax=150,
             ramp=90,  no_load=300,
             bid=45, startup=800),
}
load_mw = {"t1":150,"t2":200,
           "t3":260,"t4":240,
           "t5":180,"t6":120}
    \end{lstlisting}

    \begin{lstlisting}
m = aml.Model("mini_scuc")
G     = m.add_set("G", gen_names)
T_set = m.add_ordered_set(
            "T", t_names)  # t1..t6
    \end{lstlisting}

    \column{0.44\textwidth}
    \begin{lstlisting}
# 决策变量
u  = m.add_var2("u",  G, T_set,
      aml.VarType.Binary)
p  = m.add_var2("p",  G, T_set,
      aml.VarType.Continuous,
      lb=0.0)
su = m.add_var2("su", G, T_set,
      aml.VarType.Binary)
sd = m.add_var2("sd", G, T_set,
      aml.VarType.Binary)
    \end{lstlisting}

    \vspace{6pt}
    \begin{tcolorbox}[mybox,title=规模]
      \footnotesize
      $2 \times 6 = 12$ 连续变量\\
      $3 \times 2 \times 6 = 36$ 二值变量\\
      $\sim 60$ 个约束
    \end{tcolorbox}
  \end{columns}
\end{frame}

\begin{frame}[fragile]{AML 建模 SCUC：目标与约束}
  \begin{columns}[T]
    \column{0.52\textwidth}
    \textbf{目标函数：}
    \begin{lstlisting}
m.minimize(aml.sum_over(G, lambda g:
  aml.sum_over(T_set, lambda t:
    d["no_load"] * u[(g,t)] +
    d["bid"]     * p[(g,t)] +
    d["startup"] * su[(g,t)])))
    \end{lstlisting}

    \vspace{4pt}
    \textbf{功率平衡约束：}
    \begin{lstlisting}
for t in t_names:
  bal = aml.sum_over(G, lambda g:
          p[(g.values[0],t)] * 1.0)
  m.add_constraint(
      bal == load_mw[t],
      f"balance_{t}")
    \end{lstlisting}

    \column{0.44\textwidth}
    \textbf{出力上下界约束：}
    \begin{lstlisting}
for g, d in gen_data.items():
  for t in t_names:
    # pmin*u <= p
    lhs = (1.0*p[(g,t)]
           - d["pmin"]*u[(g,t)])
    m.add_constraint(
        lhs >= 0.0, f"lb_{g}_{t}")
    # p <= pmax*u
    rhs = (d["pmax"]*u[(g,t)]
           - 1.0*p[(g,t)])
    m.add_constraint(
        rhs >= 0.0, f"ub_{g}_{t}")
    \end{lstlisting}

    \vspace{4pt}
    \textbf{完整示例：} \texttt{demo\_04\_aml\_scuc.py}
  \end{columns}
\end{frame}

\begin{frame}[fragile]{AML 建模 SCUC：结果解读}
  \textbf{求解并提取结果：}

  \begin{columns}[T]
    \column{0.50\textwidth}
    \begin{lstlisting}
opts = aml.SolveOptions()
opts.mip_gap_tol = 0.001
result = m.solve(opts)

print(result.termination_status)
# TerminationStatus.Optimal

print(f"总费用: ${result.objective_value:,.2f}")
# 总费用: $39,250.00
    \end{lstlisting}

    \column{0.46\textwidth}
    \begin{lstlisting}
# 开停状态表
for g in gen_names:
  row = [
    f"{result.var_value(u[(g,t)]):.0f}"
    for t in t_names
  ]
  print(f"{g}: {row}")

# G1: ['1','1','1','1','1','1']
# G2: ['0','0','1','1','1','0']
    \end{lstlisting}
  \end{columns}

  \vspace{6pt}
  \textbf{典型输出：}
  \begin{center}
  \footnotesize
  \begin{tabular}{@{}lrrrrrr@{}}
    \toprule
    & t1 & t2 & t3 & t4 & t5 & t6 \\
    \midrule
    负荷 (MW) & 150 & 200 & 260 & 240 & 180 & 120 \\
    G1 状态   & 1 & 1 & 1 & 1 & 1 & 1 \\
    G1 出力 (MW) & 150.0 & 140.0 & 200.0 & 200.0 & 150.0 & 120.0 \\
    G2 状态   & 0 & 1 & 1 & 1 & 1 & 0 \\
    G2 出力 (MW) & 0.0 & 60.0 & 60.0 & 40.0 & 30.0 & 0.0 \\
    \bottomrule
  \end{tabular}
  \end{center}
\end{frame}

\begin{frame}[fragile]{JSON API：solve\_json 接口}
  \begin{columns}[T]
    \column{0.52\textwidth}
    \textbf{输入数据结构：}
    \begin{lstlisting}[language={}]
{
  "config": {
    "num_periods": 24,
    "period_length_hr": 1.0,
    "solver": "Auto",
    "mip_gap": 0.001,
    "solve_sced": true,
    "solve_lmp": true
  },
  "generators": [
    { "name": "G1", "bus": 0,
      "pmin": 50.0, "pmax": 400.0,
      "bid_segments": [
        {"price":20.0,"quantity":150.0},
        {"price":35.0,"quantity":150.0}
      ], ... }
  ],
  "loads":    [{"bus":0,"p_mw":300.0}],
  "branches": [],
  "profiles": {"load":[[0.8,...,0.95]]}
}
    \end{lstlisting}

    \column{0.44\textwidth}
    \textbf{调用与解析：}
    \begin{lstlisting}
import json, mipsolvers

json_in  = json.dumps(input_data)
json_out = mipsolvers.scuc.solve_json(
               json_in, indent=2)
res = json.loads(json_out)

scuc = res["scuc"]
print(scuc["converged"])     # True
print(scuc["cost"]["total"]) # $
    \end{lstlisting}

    \vspace{4pt}
    \textbf{输出结构：}
    \begin{lstlisting}[language={}]
{
  "scuc": {
    "converged": true,
    "cost": {
      "total": 85230.0,
      "startup": 7000.0,
      "energy": 78230.0
    },
    "dispatch": [[...],[...]],
    "commitment": [[...],[...]],
    "lmp": [20.0, 35.0, ...]
  }
}
    \end{lstlisting}
  \end{columns}
\end{frame}

\begin{frame}[fragile]{JSON API：结果分析示例}
  \begin{columns}[T]
    \column{0.52\textwidth}
    \begin{lstlisting}
dispatch   = scuc["dispatch"]
commitment = scuc["commitment"]
lmp        = scuc["lmp"]
gen_names  = ["G1", "G2", "G3"]

# 逐时段打印调度计划
print(f"{'t':>3} {'Load':>7}", end="")
for gn in gen_names:
    print(f" {gn:>8}", end="")
print(f" {'LMP':>8}")

for t in range(24):
    load_t = load_profile[t]*300.0
    print(f"{t+1:3} {load_t:7.1f}",end="")
    for gi in range(len(gen_names)):
        print(f" {dispatch[gi][t]:8.1f}",
              end="")
    print(f" {lmp[t]:8.2f}")
    \end{lstlisting}

    \column{0.44\textwidth}
    \textbf{边际电价 LMP 分析：}
    \begin{lstlisting}
import matplotlib.pyplot as plt

hours = list(range(1, 25))
fig, (ax1, ax2) = plt.subplots(
    2, 1, figsize=(8, 5))

ax1.plot(hours,
    [load_profile[t]*300 for t in range(24)])
ax1.set_ylabel("Load (MW)")

ax2.plot(hours, lmp,
         color='red', marker='o',
         markersize=3)
ax2.set_ylabel("LMP ($/MWh)")
ax2.set_xlabel("Hour")

plt.tight_layout()
plt.savefig("lmp_curve.png")
    \end{lstlisting}

    \vspace{4pt}
    \textbf{完整示例：} \texttt{demo\_05\_scuc\_api.py}
  \end{columns}
\end{frame}

% ============================================================
\section{工作进展与计划}
% ============================================================

\begin{frame}{已完成工作}
  \begin{columns}[T]
    \column{0.50\textwidth}
    \textbf{核心框架：}
    \begin{itemize}
      \item[$\checkmark$] Engine 层：LP/QP/NLP/MILP/MINLP/LE/NLE 统一接口
      \item[$\checkmark$] AML 代数建模库（16头文件 + 2实现文件）
      \item[$\checkmark$] HiGHS 集成（LP 及 MILP 开源求解器）
      \item[$\checkmark$] SCIP 集成（大规模 MILP 开源求解器）
      \item[$\checkmark$] 本地自研求解器（内点法、PDLP、LCQP、分支定界）
    \end{itemize}

    \column{0.46\textwidth}
    \textbf{Python 接口：}
    \begin{itemize}
      \item[$\checkmark$] pybind11 绑定，Python 3.8 兼容
      \item[$\checkmark$] \texttt{mipsolvers.engine} 底层接口
      \item[$\checkmark$] \texttt{mipsolvers.aml} 代数建模接口
      \item[$\checkmark$] \texttt{mipsolvers.scuc} SCUC专用接口
    \end{itemize}

    \vspace{6pt}
    \textbf{文档：}
    \begin{itemize}
      \item[$\checkmark$] Python API 参考手册（\texttt{docs/aml\_python\_api.md}）
      \item[$\checkmark$] 5个可运行 Python 示例
      \item[$\checkmark$] 本教程幻灯片
    \end{itemize}
  \end{columns}
\end{frame}

\begin{frame}{近期工作计划}
  \begin{columns}[T]
    \column{0.50\textwidth}
    \textbf{1. AC 最优潮流（AC-OPF）}
    \begin{itemize}
      \item 基于 NLP 接口实现完整 AC-OPF
      \item 支持 IEEE 测试系统（14/30/118 节点）
      \item SDP 松弛近似（凸化）
    \end{itemize}

    \vspace{8pt}
    \textbf{2. Benders 分解算法}
    \begin{itemize}
      \item 支持大规模 SCUC 加速求解
      \item 主问题（Unit Commitment）+ 子问题（SCED）解耦
      \item 适用于 $G > 200$ 台机组的场景
    \end{itemize}

    \vspace{8pt}
    \textbf{3. 验证测试集}
    \begin{itemize}
      \item IEEE RTS-96 标准测试系统
      \item 对比 Gurobi / HiGHS 求解质量与速度
    \end{itemize}

    \column{0.46\textwidth}
    \textbf{4. 网络安全约束（N-1）}
    \begin{itemize}
      \item 将输电网络约束（PTDF）纳入 SCUC
      \item N-1 准则：单一故障后仍满足约束
    \end{itemize}

    \vspace{8pt}
    \textbf{5. 分段报价完整支持}
    \begin{itemize}
      \item SCUC 模块支持多段线性报价曲线
      \item 边际电价（LMP）精确计算
      \item 竞价区域划分
    \end{itemize}

    \vspace{8pt}
    \begin{tcolorbox}[mybox,title=里程碑]
      \footnotesize
      \textbf{目标：}在标准测试系统上达到与商业求解器相当的精度，求解时间 $<5$ 分钟
    \end{tcolorbox}
  \end{columns}
\end{frame}

\begin{frame}{中期规划}
  \begin{columns}[T]
    \column{0.50\textwidth}
    \textbf{随机机组组合（Stochastic UC）}
    \begin{itemize}
      \item 考虑风电、光伏出力不确定性
      \item 基于情景树（Scenario Tree）
      \item 两阶段随机规划框架：
        \begin{itemize}
          \item 第一阶段：机组开停决策（here-and-now）
          \item 第二阶段：实时调度调整（wait-and-see）
        \end{itemize}
    \end{itemize}

    \vspace{8pt}
    \textbf{鲁棒机组组合（Robust UC）}
    \begin{itemize}
      \item 不确定性用集合描述（椭球体、多面体）
      \item 列约束生成（C\&CG）算法
    \end{itemize}

    \column{0.46\textwidth}
    \textbf{多时间尺度协同}
    \begin{itemize}
      \item 日前（24h）$\times$ 日内滚动（4h）$\times$ 实时（15min）
      \item 储能（电池）的多时段套利建模
      \item 需求响应（DR）资源聚合
    \end{itemize}

    \vspace{8pt}
    \textbf{接口扩展}
    \begin{itemize}
      \item Julia / MATLAB 绑定
      \item REST API 服务化部署
      \item 与 PowerModels.jl 对接
    \end{itemize}

    \vspace{8pt}
    \begin{tcolorbox}[mybox,title=长期目标]
      \footnotesize
      打造电力系统运行优化全栈解决方案，覆盖从稳态分析到市场出清的完整链路
    \end{tcolorbox}
  \end{columns}
\end{frame}

\begin{frame}{总结}
  \begin{columns}[T]
    \column{0.55\textwidth}
    \textbf{MIPSolvers 提供了：}
    \begin{enumerate}
      \item \textbf{统一求解器接口} — LP / MILP / NLP / MINLP / LE / NLE，通过 \texttt{solver="Auto"} 自动调度 HiGHS / SCIP / 自研求解器
      \item \textbf{代数建模库 (AML)} — 以接近数学公式的方式建模，支持多维索引变量和批量约束生成
      \item \textbf{SCUC 专用模块} — JSON 输入/输出，自动完成 UC + SCED + LMP 计算
      \item \textbf{Python 3.8 绑定} — 与 NumPy/Pandas/Matplotlib 无缝配合
    \end{enumerate}

    \column{0.40\textwidth}
    \textbf{配套示例文件：}
    \begin{itemize}
      \footnotesize
      \item \texttt{demo\_01\_lp\_direct.py}\\生产计划 LP
      \item \texttt{demo\_02\_milp\_knapsack.py}\\0/1 背包 MILP
      \item \texttt{demo\_03\_aml\_transport.py}\\运输问题（AML）
      \item \texttt{demo\_04\_aml\_scuc.py}\\迷你 SCUC（AML）
      \item \texttt{demo\_05\_scuc\_api.py}\\完整 SCUC（JSON API）
    \end{itemize}

    \vspace{8pt}
    \begin{tcolorbox}[mybox,title=运行方式]
      \footnotesize
      \texttt{python3.8 demo\_XX.py}
    \end{tcolorbox}
  \end{columns}

  \vspace{16pt}
  \begin{center}
    \large\textbf{谢谢！}
  \end{center}
\end{frame}

% ============================================================
\appendix
% ============================================================

\begin{frame}[fragile]{附录：完整安装与构建流程}
  \begin{lstlisting}[language=bash]
# 1. 克隆仓库
git clone https://github.com/.../MIPSolvers.git
cd MIPSolvers

# 2. 构建 C++ 库与测试
cmake -B build_mipsolvers -DCMAKE_BUILD_TYPE=Release
cmake --build build_mipsolvers --parallel 8
cd build_mipsolvers && ctest --output-on-failure

# 3. 构建 Python 模块 (.so)
cmake -B build_py \
      -DCMAKE_BUILD_TYPE=Release \
      -DPython_EXECUTABLE=$(which python3.8)
cmake --build build_py --parallel 8

# 4. 测试 Python 模块
python3.8 -c "
import sys; sys.path.insert(0, 'build_py')
import mipsolvers
print(mipsolvers.engine.list_solvers('LP'))
print(mipsolvers.engine.list_solvers('MILP'))
"
  \end{lstlisting}
\end{frame}

\begin{frame}{附录：求解器选择指南}
  \footnotesize
  \begin{center}
  \begin{tabular}{@{}lllll@{}}
    \toprule
    \textbf{问题类型} & \textbf{规模} & \textbf{首选求解器} & \textbf{备选} & \textbf{说明} \\
    \midrule
    LP  & 小/中 & HiGHS & NativeLCQP  & 开源，速度快 \\
    LP  & 大型 & Gurobi & HiGHS       & 商业许可证 \\
    LP  & 流式/在线 & NativePDLP & ---  & 一阶方法，低内存 \\
    MILP & 小/中 & HiGHS & SCIP        & 开源首选 \\
    MILP & 大型 & Gurobi & SCIP        & 商业许可证 \\
    MILP & 研究 & SCIP   & NativeBnC   & 强切平面支持 \\
    QP  & 凸   & HiGHS  & NativeLCQP  & 需 $Q\succeq 0$ \\
    NLP & 稀疏 & NativeIPMLP & ---     & 内点法 \\
    MINLP & 通用 & NativeBnC & ---     & 自研分支定界 \\
    \bottomrule
  \end{tabular}
  \end{center}

  \vspace{8pt}
  \textbf{注意事项：}
  \begin{itemize}
    \item \texttt{solver="Auto"} 会按优先级（Gurobi $>$ HiGHS $>$ SCIP $>$ Native...）自动选择
    \item Gurobi 需要有效许可证，否则 Auto 会跳过
    \item 对于 SCUC（$G<100,T=24$），HiGHS 通常可在 1 分钟内求解
  \end{itemize}
\end{frame}

\end{document}
